% ============================================================
% §18 — Kit Catalog & Distribution
% Author: Leslie (LaTeX Specialist)
% Last updated: 2026-04-21
% ============================================================

\chapter{Kit Catalog \& Distribution}
\label{sec:kit-catalog}

The gert v2 domain kit model (Chapter~\ref{ch:domain-kits}) establishes the architecture for domain-specific authoring semantics and tool definitions. The mobile execution model (Chapter~\ref{ch:mobile-execution}) defines the kit bundle format and platform-specific tool dispatch. This chapter specifies the mechanism by which developers \emph{discover, select, and fetch} domain kits for their applications.

Without a discovery mechanism, developers must know the exact GitHub repository URL of every kit they wish to use---a poor developer experience that does not scale. The Kit Catalog \& Distribution system solves this by providing a centralised, machine-readable index of available kits and a CLI workflow for dependency management.

% ============================================================
\section{Overview and Motivation}
\label{sec:catalog-overview}
% ============================================================

The kit catalog system is modeled after proven package distribution patterns: Homebrew formulae, CocoaPods specs repositories, and npm registries. The design principle is simple: a single, canonical repository hosts a YAML index of all available kits, and the gert CLI provides commands to search, add, and fetch kits from that index.

\subsection{Design Goals}

The catalog system satisfies three requirements:

\begin{enumerate}
  \item \textbf{Discoverability:} Developers can search for kits by name, description, or tags without prior knowledge of repository locations.
  \item \textbf{Declarative dependencies:} Applications declare their kit dependencies in a version-controlled manifest file (\texttt{Kitfile.yaml}), enabling reproducible builds and dependency auditing.
  \item \textbf{Decentralised publishing:} Kit authors publish to the catalog via pull requests to a public repository. There is no central gatekeeper, approval process, or proprietary registry service.
\end{enumerate}

\subsection{Non-Goals}

The catalog is explicitly \emph{not} a package hosting service. It does not store kit bundles, serve downloadable artifacts, or maintain versioned releases. The catalog is purely a discovery index---a pointer to source repositories. Kit bundles are fetched directly from their declared source locations (typically GitHub repositories).

% ============================================================
\section{The Catalog Repository}
\label{sec:catalog-repo}
% ============================================================

The canonical kit catalog is maintained in the GitHub repository:

\begin{center}
\texttt{ormasoftchile/gert-catalog}
\end{center}

This repository contains a single file: \texttt{catalog.yaml}, which serves as the machine-readable index of all published kits.

\subsection{Repository Structure}

\begin{verbatim}
gert-catalog/
  catalog.yaml        # The catalog index
  README.md           # Publishing guidelines for kit authors
  LICENSE             # MIT License
\end{verbatim}

The \texttt{README.md} provides instructions for kit authors on how to submit new entries. The catalog itself is fully self-describing---every field in \texttt{catalog.yaml} is documented inline via YAML comments.

\subsection{Access and Availability}

The catalog repository is public and requires no authentication to read. The \texttt{catalog.yaml} file is served via GitHub's raw content CDN:

\begin{center}
\texttt{https://raw.githubusercontent.com/ormasoftchile/gert-catalog/main/catalog.yaml}
\end{center}

This URL is the default catalog source for all gert CLI installations. Organisations may fork the catalog repository and configure the CLI to reference their internal fork for private kit distribution.

% ============================================================
\section{Catalog Schema}
\label{sec:catalog-schema}
% ============================================================

The \texttt{catalog.yaml} file is a flat list of kit entries, each describing a single published kit.

\subsection{Schema Definition}

\begin{minted}{yaml}
# catalog.yaml — Kit Catalog Index
apiVersion: catalog/v1

kits:
  - name: gert-mobile-platform
    description: >
      Camera, location, NFC, biometrics, and notifications 
      for mobile runbooks
    source: github:ormasoftchile/gert-mobile-platform
    latest: "0.1.0"
    targets: [ios, android]
    tags: [mobile, platform, capabilities]

  - name: home
    description: Home automation runbooks and tools
    source: github:ormasoftchile/gert-domain-home
    latest: "1.0.0"
    targets: [ios, android, cli]
    tags: [home, iot, smart-home]

  - name: pool-maintenance
    description: Pool maintenance inspection and reporting runbooks
    source: github:ormasoftchile/gert-domain-pool
    latest: "1.0.0"
    targets: [ios, android]
    tags: [facilities, maintenance]
\end{minted}

\subsection{Field Descriptions}

\begin{description}
  \item[\texttt{name}] (required) — The unique identifier for the kit. This is the name used in \texttt{Kitfile.yaml} and in CLI commands (\texttt{gert kit add home}). By convention, platform kits use the \texttt{gert-} prefix; domain kits do not.

  \item[\texttt{description}] (required) — A human-readable summary of the kit's purpose. This appears in search results and kit listing output.

  \item[\texttt{source}] (required) — The location from which the kit bundle can be fetched. The format is \texttt{provider:path}, where \texttt{provider} is one of:
  \begin{itemize}
    \item \texttt{github:owner/repo} — Fetch from a GitHub repository
    \item \texttt{git:url} — Fetch from an arbitrary Git repository URL
    \item \texttt{https:url} — Fetch from an HTTPS endpoint (expects a \texttt{.kit} archive)
  \end{itemize}
  Currently, only GitHub sources are fully implemented.

  \item[\texttt{latest}] (required) — The most recent stable version of the kit, expressed as a semantic version string. This is used when a \texttt{Kitfile.yaml} declares a dependency without specifying a version constraint.

  \item[\texttt{targets}] (optional) — An array of supported execution platforms. Valid values are \texttt{ios}, \texttt{android}, \texttt{cli}, and \texttt{server}. If omitted, the kit is assumed to support all platforms.

  \item[\texttt{tags}] (optional) — An array of keywords for categorisation and search filtering. Tags are freeform but should follow a consistent taxonomy (e.g., \texttt{mobile}, \texttt{iot}, \texttt{compliance}, \texttt{workflow}).
\end{description}

\subsection{Validation Rules}

Kit entries in \texttt{catalog.yaml} are subject to the following validation rules, enforced by the catalog maintainer during pull request review:

\begin{itemize}
  \item Kit names must be unique within the catalog.
  \item Kit names must match the pattern \texttt{[a-z0-9-]+} (lowercase alphanumeric and hyphens only).
  \item The \texttt{source} field must reference a publicly accessible repository or endpoint.
  \item The \texttt{latest} field must be a valid semantic version string (\texttt{MAJOR.MINOR.PATCH}).
  \item Tags should be lowercase and hyphen-separated.
\end{itemize}

% ============================================================
\section{Application Kit Dependencies: Kitfile.yaml}
\label{sec:kitfile}
% ============================================================

An application project declares its kit dependencies in a \texttt{Kitfile.yaml} file at the project root. This file serves the same purpose as \texttt{package.json} (npm), \texttt{Gemfile} (Ruby), or \texttt{requirements.txt} (Python)---it is a declarative manifest of external dependencies.

\subsection{Schema and Example}

\begin{minted}{yaml}
# Kitfile.yaml — Application Kit Dependencies
catalog: https://raw.githubusercontent.com/ormasoftchile/gert-catalog/main/catalog.yaml

kits:
  - name: home
    version: "^1.0.0"
  - name: gert-mobile-platform
    version: "^0.1.0"
\end{minted}

\subsection{Field Descriptions}

\begin{description}
  \item[\texttt{catalog}] (required) — The URL of the catalog index to use for kit resolution. The default value is the canonical \texttt{gert-catalog} repository. Organisations may override this to reference a private fork or internal registry.

  \item[\texttt{kits}] (required) — An array of kit dependency entries, each containing:
  \begin{itemize}
    \item \texttt{name} (required) — The kit identifier, matching a \texttt{name} field in \texttt{catalog.yaml}.
    \item \texttt{version} (optional) — A semantic version constraint (see \S\ref{sec:version-resolution}). If omitted, the \texttt{latest} version from the catalog is used.
  \end{itemize}
\end{description}

\subsection{Version Constraints}

Version constraints follow the npm semver syntax:

\begin{table}[ht]
\centering
\caption{Semantic version constraint syntax}
\label{tab:semver-constraints}
\small
\begin{tabular}{p{0.25\linewidth}p{0.55\linewidth}}
\toprule
\textbf{Constraint} & \textbf{Meaning} \\
\midrule
\texttt{\^{}1.0.0} & Compatible with 1.0.0 (allows minor and patch updates: \texttt{>=1.0.0 <2.0.0}) \\
\texttt{\~{}1.2.0} & Approximately 1.2.0 (allows patch updates only: \texttt{>=1.2.0 <1.3.0}) \\
\texttt{>=1.0.0 <2.0.0} & Explicit range (any version satisfying both constraints) \\
\texttt{1.0.0} & Exact version (no flexibility) \\
\texttt{*} & Any version (not recommended for production) \\
\bottomrule
\end{tabular}
\end{table}

The most commonly used constraint is \texttt{\^{}MAJOR.MINOR.PATCH} (caret range), which allows backward-compatible updates within the same major version.

% ============================================================
\section{CLI Commands}
\label{sec:catalog-cli}
% ============================================================

The gert CLI provides a \texttt{gert kit} subcommand with four operations: search, add, fetch, and list.

\subsection{Command Reference}

\begin{table}[ht]
\centering
\caption{gert kit subcommands}
\label{tab:kit-commands}
\small
\begin{tabular}{p{0.30\linewidth}p{0.55\linewidth}}
\toprule
\textbf{Command} & \textbf{Description} \\
\midrule
\texttt{gert kit search <query>} &
  Fetches \texttt{catalog.yaml} from the configured catalog source, filters entries by name, description, or tags matching \texttt{<query>}, and prints results in a tabular format. \\
\midrule
\texttt{gert kit add <name> [\texttt{-{}-version} <constraint>]} &
  Adds a kit dependency to \texttt{Kitfile.yaml}. If \texttt{-{}-version} is omitted, uses \texttt{latest} from the catalog. Creates \texttt{Kitfile.yaml} if it does not exist. \\
\midrule
\texttt{gert kit fetch [\texttt{-{}-target} <platform>]} &
  Resolves all kits declared in \texttt{Kitfile.yaml}, downloads \texttt{.kit} bundles to \texttt{.gert/kits/}, resolves transitive dependencies, and generates \texttt{Kitfile.lock}. If \texttt{-{}-target} is specified, only fetches kits compatible with that platform. \\
\midrule
\texttt{gert kit list} &
  Displays all kits currently fetched in \texttt{.gert/kits/}, showing name, version, and source. \\
\bottomrule
\end{tabular}
\end{table}

\subsection{Command Examples}

\paragraph{Searching for kits.}

\begin{verbatim}
$ gert kit search home
NAME                DESCRIPTION                              LATEST   TARGETS
home                Home automation runbooks and tools       1.0.0    ios, android, cli
pool-maintenance    Pool maintenance inspection runbooks     1.0.0    ios, android
\end{verbatim}

\paragraph{Adding a kit dependency.}

\begin{verbatim}
$ gert kit add home --version "^1.0.0"
✓ Added kit 'home' to Kitfile.yaml with version constraint '^1.0.0'
\end{verbatim}

\paragraph{Fetching all kits.}

\begin{verbatim}
$ gert kit fetch
Resolving dependencies...
  ✓ home@1.0.0
  ✓ gert-mobile-platform@0.1.0 (transitive dependency of home)
Downloading kits...
  ✓ home@1.0.0 → .gert/kits/home-1.0.0.kit/
  ✓ gert-mobile-platform@0.1.0 → .gert/kits/gert-mobile-platform-0.1.0.kit/
Generated Kitfile.lock
\end{verbatim}

\paragraph{Listing fetched kits.}

\begin{verbatim}
$ gert kit list
NAME                    VERSION   SOURCE
home                    1.0.0     github:ormasoftchile/gert-domain-home
gert-mobile-platform    0.1.0     github:ormasoftchile/gert-mobile-platform
\end{verbatim}

% ============================================================
\section{Publishing a Kit}
\label{sec:kit-publishing}
% ============================================================

Publishing a kit to the catalog is a manual, pull-request-based workflow. This design is intentional: it provides a lightweight review gate without requiring a centralised authority or approval process.

\subsection{Publishing Workflow}

\begin{enumerate}
  \item \textbf{Create the kit repository.} The kit author creates a GitHub repository (e.g., \texttt{ormasoftchile/gert-domain-weather}) containing the kit's source code, manifest, compiler, and tool definitions. The repository must include a \texttt{manifest.json} or \texttt{gert-kit.yaml} file at the root.

  \item \textbf{Tag a release.} The author creates a Git tag matching the kit version (e.g., \texttt{v1.0.0}). This tag is used by the \texttt{gert kit fetch} command to download a specific version.

  \item \textbf{Fork the catalog repository.} The author forks \texttt{ormasoftchile/gert-catalog} to their own GitHub account.

  \item \textbf{Add a catalog entry.} The author edits \texttt{catalog.yaml} to add a new entry for their kit, following the schema defined in \S\ref{sec:catalog-schema}.

  \item \textbf{Open a pull request.} The author submits a pull request from their fork to the upstream \texttt{gert-catalog} repository. The pull request description should include:
  \begin{itemize}
    \item A brief summary of the kit's purpose.
    \item A link to the kit's repository.
    \item Confirmation that the repository is publicly accessible.
  \end{itemize}

  \item \textbf{Review and merge.} The catalog maintainer reviews the pull request to ensure:
  \begin{itemize}
    \item The kit name is unique.
    \item The \texttt{source} URL is valid and publicly accessible.
    \item The \texttt{latest} version tag exists in the source repository.
    \item The entry follows the catalog schema.
  \end{itemize}
  Once approved, the maintainer merges the pull request, making the kit immediately discoverable via \texttt{gert kit search}.
\end{enumerate}

\subsection{Updating a Published Kit}

When a kit author releases a new version, they submit a pull request to update the \texttt{latest} field in \texttt{catalog.yaml}. The maintainer verifies that the new version tag exists in the source repository and merges the update.

Kit entries are never removed from the catalog unless the source repository is deleted or the kit is deprecated. Deprecated kits may be marked with a \texttt{deprecated: true} field (future extension to the catalog schema).

% ============================================================
\section{Multi-Kit Bundling}
\label{sec:multi-kit}
% ============================================================

A single \texttt{Kitfile.yaml} may declare dependencies on multiple kits. The \texttt{gert kit fetch} command resolves all declared kits and downloads them to \texttt{.gert/kits/}.

\subsection{Capability Collision Resolution}

When multiple kits provide implementations for the same tool (e.g., both \texttt{home} and \texttt{pool-maintenance} define a \texttt{camera.capture} tool), the \textbf{last-write-wins} strategy is used. The order of kits in \texttt{Kitfile.yaml} determines precedence:

\begin{minted}{yaml}
kits:
  - name: home
    version: "^1.0.0"
  - name: pool-maintenance
    version: "^1.0.0"
\end{minted}

In this example, if both \texttt{home} and \texttt{pool-maintenance} define \texttt{camera.capture}, the implementation from \texttt{pool-maintenance} is used because it appears later in the list.

This order-dependent behaviour is explicit and documented. Developers are encouraged to avoid capability collisions by using namespaced tool names (e.g., \texttt{home.camera.capture} vs. \texttt{pool.camera.capture}) when authoring domain kits.

\subsection{Tool Registry Merge}

When \texttt{gert kit fetch} completes, the CLI merges all tool definitions from all fetched kits into a unified tool registry at \texttt{.gert/tools/}. This registry is queried at runtime by the executor when a runbook step references a tool.

% ============================================================
\section{Platform-Aware Kits}
\label{sec:platform-aware}
% ============================================================

Kit entries in \texttt{catalog.yaml} include an optional \texttt{targets} field that declares the platforms the kit supports. This allows the \texttt{gert kit fetch} command to filter kits based on the target platform.

\subsection{Platform-Filtered Fetching}

When invoked with the \texttt{-{}-target} flag, \texttt{gert kit fetch} only downloads kits that declare compatibility with the specified platform:

\begin{verbatim}
$ gert kit fetch --target ios
Resolving dependencies...
  ✓ home@1.0.0 (targets: ios, android, cli)
  ✓ gert-mobile-platform@0.1.0 (targets: ios, android)
Skipping 'server-monitoring' (targets: server, cli)
\end{verbatim}

This is particularly useful for mobile application builds, where server-only kits should be excluded.

\subsection{Platform-Specific Tool Implementations}

As described in Chapter~\ref{ch:mobile-execution}, individual tool YAML files within a kit encode platform-specific handlers via the \texttt{impl} block:

\begin{minted}{yaml}
impl:
  ios:
    transport: native-sdk
    handler: GertSDK.Camera.capture
  android:
    transport: native-sdk
    handler: com.gert.platform.tools.CameraCaptureTool
  cli:
    transport: binary
    handler: ./bin/camera-capture-cli
\end{minted}

This allows a single kit to provide cross-platform tool definitions while delegating platform-specific behaviour to appropriate handlers. The \texttt{targets} field in the catalog entry reflects the union of all platforms for which at least one tool implementation exists.

% ============================================================
\section{Kit-to-Kit Dependencies}
\label{sec:kit-dependencies}
% ============================================================

A domain kit may depend on one or more other kits, typically platform kits. Kit-to-kit dependencies are declared in the kit's \texttt{manifest.json} file via the \texttt{requires} field.

\subsection{Dependency Declaration}

\begin{minted}{json}
{
  "name": "home",
  "version": "1.0.0",
  "requires": [
    {
      "kit": "gert-mobile-platform",
      "version": "^0.1.0"
    }
  ]
}
\end{minted}

The \texttt{requires} field is an array of dependency objects, each containing:

\begin{itemize}
  \item \texttt{kit} — The name of the required kit (matching a \texttt{name} in \texttt{catalog.yaml}).
  \item \texttt{version} — A semantic version constraint.
\end{itemize}

\subsection{Transitive Dependency Resolution}

When \texttt{gert kit fetch} processes a kit, it recursively resolves all \texttt{requires} entries:

\begin{enumerate}
  \item Fetch the top-level kit declared in \texttt{Kitfile.yaml}.
  \item Parse the kit's \texttt{manifest.json}.
  \item For each entry in \texttt{requires}, resolve the required kit from the catalog.
  \item Recursively fetch and resolve dependencies of the required kit.
  \item Deduplicate kits by name, selecting the highest compatible version when multiple constraints overlap.
\end{enumerate}

This is a depth-first resolution strategy identical to npm and Bundler. Circular dependencies are detected and result in a resolution error.

\subsection{Dependency Graph Example}

\begin{verbatim}
Application (Kitfile.yaml)
  ├── home@1.0.0
  │     └── gert-mobile-platform@0.1.0
  └── pool-maintenance@1.0.0
        └── gert-mobile-platform@0.1.0

Resolved:
  - home@1.0.0
  - pool-maintenance@1.0.0
  - gert-mobile-platform@0.1.0  (shared dependency, resolved once)
\end{verbatim}

% ============================================================
\section{Version Resolution and Locking}
\label{sec:version-resolution}
% ============================================================

Version resolution follows semantic versioning principles and produces a deterministic, reproducible set of kit versions for a given \texttt{Kitfile.yaml}.

\subsection{Resolution Algorithm}

The \texttt{gert kit fetch} command performs version resolution in two passes:

\begin{enumerate}
  \item \textbf{Constraint collection:} Build a dependency graph by recursively resolving \texttt{requires} entries. Collect all version constraints for each kit name.
  
  \item \textbf{Constraint solving:} For each kit, compute the intersection of all collected constraints. Select the highest version from \texttt{catalog.yaml} that satisfies the intersection. If no version satisfies all constraints, emit a resolution error.
\end{enumerate}

\subsection{Example Resolution}

Given the following constraints:

\begin{itemize}
  \item \texttt{home@1.0.0} requires \texttt{gert-mobile-platform@\^{}0.1.0} (i.e., \texttt{>=0.1.0 <1.0.0})
  \item \texttt{pool-maintenance@1.0.0} requires \texttt{gert-mobile-platform@\~{}0.1.0} (i.e., \texttt{>=0.1.0 <0.2.0})
\end{itemize}

The constraint intersection is \texttt{>=0.1.0 <0.2.0}. If the catalog lists \texttt{gert-mobile-platform} versions \texttt{0.1.0}, \texttt{0.1.5}, and \texttt{0.2.0}, the resolver selects \texttt{0.1.5} (the highest version within the constraint range).

\subsection{Kitfile.lock}

After successful resolution, \texttt{gert kit fetch} generates a \texttt{Kitfile.lock} file that records the exact resolved versions:

\begin{minted}{yaml}
# Kitfile.lock — Generated by gert kit fetch
# Do not edit manually. Commit to version control for reproducible builds.

resolved:
  - name: home
    version: "1.0.0"
    source: github:ormasoftchile/gert-domain-home
    checksum: a3f9c1...

  - name: pool-maintenance
    version: "1.0.0"
    source: github:ormasoftchile/gert-domain-pool
    checksum: b2e4d5...

  - name: gert-mobile-platform
    version: "0.1.5"
    source: github:ormasoftchile/gert-mobile-platform
    checksum: c7d8e9...
\end{minted}

The lockfile includes:

\begin{itemize}
  \item \texttt{name} — The kit identifier.
  \item \texttt{version} — The exact resolved version.
  \item \texttt{source} — The source location from which the kit was fetched.
  \item \texttt{checksum} — A SHA256 hash of the fetched kit bundle, ensuring integrity across installations.
\end{itemize}

\subsection{Lockfile Workflow}

The lockfile should be committed to version control alongside \texttt{Kitfile.yaml}. This ensures that all team members and CI/CD pipelines fetch identical kit versions, preventing ``works on my machine'' discrepancies.

When \texttt{Kitfile.lock} is present, \texttt{gert kit fetch} uses the locked versions instead of re-resolving constraints. To update dependencies to their latest compatible versions, delete \texttt{Kitfile.lock} and re-run \texttt{gert kit fetch}.

% ============================================================
\section{Relationship to Existing Design}
\label{sec:catalog-relationships}
% ============================================================

The kit catalog system builds upon the foundations established in earlier chapters:

\begin{description}
  \item[Domain Kit Model (Chapter~\ref{ch:domain-kits})] defines the architecture and lifecycle of domain kits. The catalog provides the \emph{discovery and distribution} layer for these kits, answering the question ``how do developers find and fetch kits?''

  \item[Mobile Execution (Chapter~\ref{ch:mobile-execution})] defines the kit bundle format (\texttt{.kit} directories) and the \texttt{manifest.json} schema. The catalog references these bundles via the \texttt{source} field and resolves \texttt{requires} dependencies declared in manifests.

  \item[Tool Runtime (Chapter~\ref{ch:tools})] defines the tool YAML schema and execution contract. Kits distributed via the catalog contain tool definitions that conform to this schema. The merged tool registry produced by \texttt{gert kit fetch} is consumed by the tool runtime at execution time.
\end{description}

The catalog is not a replacement for any of these systems---it is the distribution and dependency management layer that connects them.

% ============================================================
\section{Alternative Catalog Sources}
\label{sec:alternative-catalogs}
% ============================================================

While the canonical catalog is \texttt{ormasoftchile/gert-catalog}, organisations may choose to host private catalogs for internal kits.

\subsection{Configuring a Custom Catalog}

The catalog URL is configurable via the \texttt{catalog} field in \texttt{Kitfile.yaml}:

\begin{minted}{yaml}
catalog: https://git.acme.example/gert/internal-catalog/raw/main/catalog.yaml

kits:
  - name: acme-enterprise-platform
    version: "^2.0.0"
\end{minted}

This allows enterprises to:

\begin{itemize}
  \item Maintain a curated set of approved kits.
  \item Distribute proprietary kits that cannot be open-sourced.
  \item Enforce corporate governance policies (e.g., all kits must be signed by the internal CA).
\end{itemize}

\subsection{Catalog Merging}

The CLI does not currently support merging multiple catalogs. If an application requires kits from both the public catalog and a private catalog, the private catalog must re-declare entries from the public catalog. This is intentional: it ensures that the private catalog is the authoritative source for all kit metadata.

% ============================================================
\section{Security Considerations}
\label{sec:catalog-security}
% ============================================================

The catalog system introduces several trust boundaries that must be carefully managed.

\subsection{Catalog Source Trust}

The \texttt{catalog} URL in \texttt{Kitfile.yaml} is a critical trust anchor. If an attacker can modify this field (e.g., via a malicious pull request), they can redirect kit fetches to attacker-controlled repositories.

\paragraph{Mitigation:} The \texttt{Kitfile.yaml} file should be version-controlled and subject to code review. Changes to the \texttt{catalog} field should be treated with the same scrutiny as changes to dependency versions.

\subsection{Kit Source Integrity}

The \texttt{source} field in \texttt{catalog.yaml} points to external repositories. If the catalog is compromised, an attacker could replace a legitimate \texttt{source} URL with a malicious one.

\paragraph{Mitigation:} The \texttt{Kitfile.lock} file includes a \texttt{checksum} field for each kit. On subsequent fetches, the CLI verifies that the downloaded kit matches the locked checksum. This prevents silent substitution attacks.

\subsection{Kit Bundle Signing}

As described in Chapter~\ref{ch:mobile-execution}, kit bundles may be signed using Ed25519 signatures. The catalog schema does not currently include a \texttt{signingKey} field, but this is a natural extension for future versions.

\paragraph{Future extension:} Add a \texttt{signingKey} field to catalog entries, referencing a public key or keyserver. The CLI would verify bundle signatures against the declared key before loading the kit.

% ============================================================
\section{Future Extensions}
\label{sec:catalog-future}
% ============================================================

The catalog system is designed to be extensible. The following features are under consideration for future versions:

\begin{description}
  \item[Deprecation markers.] Add a \texttt{deprecated: true} field to catalog entries, allowing kit authors to signal that a kit is no longer maintained. The CLI would emit a warning when fetching deprecated kits.

  \item[Platform-specific versions.] Allow catalog entries to declare different \texttt{latest} versions for different platforms (e.g., \texttt{latest\_ios: "1.2.0"}, \texttt{latest\_android: "1.1.5"}). This accommodates platform-specific release cadences.

  \item[Alias support.] Allow catalog entries to declare aliases (e.g., \texttt{home-automation} as an alias for \texttt{home}). This improves discoverability without requiring kit renames.

  \item[Catalog versioning.} Introduce a \texttt{catalog/v2} schema with structured metadata (e.g., author contact information, homepage URL, license type). The CLI would support both \texttt{catalog/v1} and \texttt{catalog/v2} during a transition period.
\end{description}

These extensions require changes to the catalog schema and CLI implementation but do not alter the fundamental architecture.
